%==============================================================================
% CAPÍTULO 16 — AUDITORIA DE REPOSITÓRIOS GITHUB
% PARTE III — Engenharia de Software Odontológico com SDD e TDD
%==============================================================================
\chapter{Auditoria de Repositórios GitHub}
\label{cap:auditoria-github}
\nivelquatro

O GitHub hospeda milhares de repositórios de IA odontológica. Alguns são
tesouros de pesquisa com código de produção, artigos publicados e modelos
pré-treinados. Outros são projetos abandonados de fim de semana, com
código quebrado e nenhuma documentação. Este capítulo ensina como
\textbf{distinguir joias de cascalhos} — e como aplicar TDD sobre
código de terceiros.

\section{Como Pesquisar IA Odontológica no GitHub}

\subsection{Estratégias de busca eficazes}

A busca padrão do GitHub é limitada. Use estas estratégias avançadas
para encontrar repositórios odontológicos de qualidade:

\begin{table}[H]
  \centering
  \caption{Estratégias de busca no GitHub para IA odontológica.}
  \label{tab:github-search}
  \begin{tabularx}{\textwidth}{lX}
    \toprule
    \textbf{Query} & \textbf{O que encontra} \\
    \midrule
    \texttt{dental deep learning stars:>50} & Repositórios populares ($>$50 stars) \\
    \texttt{dental segmentation language:Python} & Código Python de segmentação dental \\
    \texttt{orthodontic AI pushed:>2024-01-01} & Atualizados recentemente \\
    \texttt{caries detection topic:deep-learning} & Com tópico ``deep-learning'' \\
    \texttt{panoramic radiograph license:mit} & Com licença MIT (reutilizável) \\
    \texttt{oral cancer classification dataset} & Com dataset incluso ou linkado \\
    \texttt{tooth segmentation in:readme} & Mencionado no README \\
    \bottomrule
  \end{tabularx}
\end{table}

\indice{GitHub} \indice{busca avançada}

\destaque{Combine múltiplos filtros para refinar resultados. Exemplo:
\texttt{dental segmentation language:Python stars:>30 license:mit}
encontra repositórios populares, em Python, com licença permissiva.}

\subsection{O que observar em um repositório}

Ao abrir um repositório, faça a \textbf{inspeção de 30 segundos}:

\begin{enumerate}
  \item \textbf{README.md}: existe? explica o que o projeto faz? tem imagens?
  \item \textbf{Último commit}: foi há 3 dias ou há 3 anos?
  \item \textbf{Stars e forks}: popularidade não é qualidade, mas é um sinal.
  \item \textbf{Issues abertas}: muitas issues sem resposta indicam projeto abandonado.
  \item \textbf{LICENSE}: sem licença = não pode usar legalmente.
  \item \textbf{requirements.txt / environment.yml}: sem dependências explícitas = não reproduzível.
\end{enumerate}

\section{Níveis G1--G3 de Confiabilidade}

Desenvolvemos um sistema de classificação em três níveis para avaliar a
confiabilidade de repositórios de IA odontológica. Inspirado nos níveis
de prontidão tecnológica (\textit{Technology Readiness Levels}) da NASA,
mas adaptado para software científico-clínico.

\indice{níveis G1-G3} \indice{confiabilidade de repositórios}

\subsection{Nível G1 — Básico (Exploratório)}

Um repositório G1 atende aos \textbf{requisitos mínimos de existência e
organização}. É adequado para exploração e aprendizado, mas \textbf{não}
para uso em pesquisa ou clínica.

\begin{table}[H]
  \centering
  \caption{Checklist G1 — Nível Básico.}
  \label{tab:g1-checklist}
  \begin{tabularx}{\textwidth}{cXc}
    \toprule
    \textbf{\#} & \textbf{Requisito} & \textbf{Peso} \\
    \midrule
    G1.1 & README.md existe e está em inglês ou português & Obrigatório \\
    G1.2 & Arquivo LICENSE presente & Obrigatório \\
    G1.3 & requirements.txt ou environment.yml & Obrigatório \\
    G1.4 & Código executável sem erros de sintaxe & Obrigatório \\
    G1.5 & Estrutura de diretórios organizada & Desejável \\
    \bottomrule
  \end{tabularx}
\end{table}

\subsection{Nível G2 — Intermediário (Pesquisa)}

Um repositório G2 atende aos requisitos G1 \textbf{mais} características
de rigor científico. É adequado para \textbf{reprodução de experimentos}
e \textbf{comparação em pesquisas}.

\begin{table}[H]
  \centering
  \caption{Checklist G2 — Nível Intermediário (requer G1 completo).}
  \label{tab:g2-checklist}
  \begin{tabularx}{\textwidth}{cXc}
    \toprule
    \textbf{\#} & \textbf{Requisito} & \textbf{Peso} \\
    \midrule
    G2.1 & Citação de artigo/DOI associado & Obrigatório \\
    G2.2 & Dataset documentado (fonte, tamanho, licença) & Obrigatório \\
    G2.3 & Scripts de treinamento reproduzíveis & Obrigatório \\
    G2.4 & Métricas reportadas (accuracy, sensibilidade, especificidade) & Obrigatório \\
    G2.5 & Testes automatizados (pytest/unittest) & Desejável \\
    \bottomrule
  \end{tabularx}
\end{table}

\subsection{Nível G3 — Avançado (Produção)}

Um repositório G3 atende aos requisitos G1 e G2 \textbf{mais}
características de software profissional. É adequado para \textbf{uso
clínico-assistido} e \textbf{integração em sistemas}.

\begin{table}[H]
  \centering
  \caption{Checklist G3 — Nível Avançado (requer G1 e G2 completos).}
  \label{tab:g3-checklist}
  \begin{tabularx}{\textwidth}{cXc}
    \toprule
    \textbf{\#} & \textbf{Requisito} & \textbf{Peso} \\
    \midrule
    G3.1 & Modelo pré-treinado disponível (Hugging Face, Zenodo) & Obrigatório \\
    G3.2 & Dockerfile para ambiente containerizado & Desejável \\
    G3.3 & CI/CD pipeline (GitHub Actions) & Desejável \\
    G3.4 & Model Card ou Data Sheet & Obrigatório \\
    G3.5 & Reprodutibilidade validada por terceiros & Desejável \\
    \bottomrule
  \end{tabularx}
\end{table}

\section{Protocolo de Auditoria}

O protocolo completo de auditoria consiste em 5 etapas executadas
sequencialmente:

\subsection{Etapa 1: Verificação de Licença}

\textbf{Sem licença explícita, o código é proprietário por padrão} —
você não pode usá-lo, modificá-lo ou redistribuí-lo legalmente, mesmo
que esteja público no GitHub.

\begin{itemize}
  \item Procure por \texttt{LICENSE}, \texttt{LICENSE.md} ou \texttt{COPYING}.
  \item Identifique o tipo (MIT, Apache 2.0, GPLv3, CC BY, etc.).
  \item Verifique compatibilidade com seu uso pretendido (ver Capítulo~17).
\end{itemize}

\subsection{Etapa 2: Verificação de Artigo e DOI}

Um repositório de IA odontológica \textbf{sério} está associado a uma
publicação científica. Verifique:

\begin{itemize}
  \item O artigo citado existe? Busque o DOI no \url{https://doi.org}.
  \item O artigo foi publicado em periódico com revisão por pares ou é
    apenas um preprint no arXiv?
  \item O código no repositório corresponde ao descrito no artigo?
  \item Os resultados reportados no artigo são reproduzíveis com o código?
\end{itemize}

\subsection{Etapa 3: Verificação do Dataset}

\begin{itemize}
  \item O dataset está disponível publicamente ou é privado?
  \item Se privado, há descrição suficiente para avaliar a qualidade?
  \item Há documentação sobre anonimização e aspectos éticos (CEP, CAAE)?
  \item O balanceamento de classes é adequado para a tarefa clínica?
\end{itemize}

\subsection{Etapa 4: Execução e Testes}

\begin{itemize}
  \item O código compila/executa sem erros na primeira tentativa?
  \item Os testes automatizados passam (\texttt{pytest --tb=short})?
  \item Há testes para os cenários de falha (\textit{edge cases})?
  \item A cobertura de testes é adequada ($>$80\%)?
\end{itemize}

\subsection{Etapa 5: Aplicação de TDD sobre Código de Terceiros}

Mesmo que o repositório original não tenha testes, \textbf{você pode —
e deve — adicionar seus próprios testes} antes de usar o código:

\begin{enumerate}
  \item \textbf{Fork} o repositório para sua conta.
  \item Crie uma branch \texttt{audit/tdd-tests}.
  \item Adicione testes para as funções que você pretende usar.
  \item Se os testes falharem, investigue antes de confiar no código.
  \item Documente suas descobertas em um \texttt{AUDIT\_REPORT.md}.
\end{enumerate}

\teste{TDD sobre código de terceiros: ao adicionar testes para a função
\texttt{preprocess\_dental\_image} de um repositório G2, descobriu-se que
ela não tratava imagens com dimensões ímpares, causando crash silencioso.
O bug foi reportado como issue e corrigido em fork.}

\section{Estudos de Caso}

\subsection{OdontoAI — Classification}

\begin{description}
  \item[URL] \url{https://github.com/odontoai/classification}
  \item[Nível] G2 — Intermediário (Pesquisa)
  \item[Pontos fortes] Artigo publicado (DOI: 10.1016/j.jdent.2023.104562),
    dataset documentado (3.400 radiografias), métricas reportadas.
  \item[Pontos fracos] Sem testes automatizados, sem CI/CD, sem Model Card.
  \item[Recomendação] Adequado para pesquisa com adição de testes TDD.
\end{description}

\subsection{DENTEX — Dental Enumeration and Diagnosis}

\begin{description}
  \item[URL] \url{https://github.com/ibrahimethemhamamci/DENTEX}
  \item[Nível] G3 — Avançado (Produção)
  \item[Pontos fortes] Artigo premiado no MICCAI 2023, Dockerfile,
    CI/CD ativo, modelo pré-treinado no Hugging Face, Model Card completo.
  \item[Pontos fracos] Dataset privado (não reprodutível sem acesso).
  \item[Recomendação] Pronto para uso clínico-assistido com adaptação
    ao dataset local.
\end{description}

\subsection{ToothSeg — Tooth Segmentation}

\begin{description}
  \item[URL] \url{https://github.com/liangjiubujiu/ToothSeg}
  \item[Nível] G1+ (Básico com alguns elementos G2)
  \item[Pontos fortes] Código organizado, artigo no arXiv, modelo
    pré-treinado disponível para download.
  \item[Pontos fracos] Dataset não documentado, sem testes, sem CI/CD,
    última atualização há 18 meses.
  \item[Recomendação] Útil para aprendizado de arquiteturas de
    segmentação, mas requer auditoria completa antes de uso em pesquisa.
\end{description}

\subsection{MeshSegNet — 3D Dental Mesh Segmentation}

\begin{description}
  \item[URL] \url{https://github.com/Tai-Hsien/MeshSegNet}
  \item[Nível] G2 — Intermediário (Pesquisa)
  \item[Pontos fortes] Artigo no IEEE TMI (DOI: 10.1109/TMI.2020.2977934),
    código bem documentado, suporte a mesh 3D.
  \item[Pontos fracos] Dependências complexas (Pytorch Geometric, trimesh),
    sem testes, dataset requer convênio institucional.
  \item[Recomendação] Excelente para pesquisa em gêmeos digitais 3D.
\end{description}

\subsection{MLUA — Machine Learning in Ultrasound Analysis}

\begin{description}
  \item[URL] \url{https://github.com/anonymous-dental/MLUA}
  \item[Nível] G1 — Básico
  \item[Pontos fortes] Código funcional, boa documentação inline.
  \item[Pontos fracos] Sem artigo, sem DOI, sem dataset, sem licença (!),
    autor anônimo — múltiplos sinais de alerta.
  \item[Recomendação] \textbf{Não utilizar.} Repositório fantasma: código
    sem autor identificável, sem licença e sem validação científica
    constitui risco jurídico e científico.
\end{description}

\subsection{DentalMIM — Dental Model Innovation}

\begin{description}
  \item[URL] \url{https://github.com/dentalmim/DentalMIM}
  \item[Nível] G2+ — Intermediário tendendo a Avançado
  \item[Pontos fortes] Artigo na Scientific Reports, licença MIT,
    pipeline completo de treinamento, dataset público (Kaggle).
  \item[Pontos fracos] Sem Docker, sem CI/CD, cobertura de testes parcial.
  \item[Recomendação] Bom ponto de partida para classificação de imagens
    odontológicas. Fork e adicione Docker + CI/CD para elevar a G3.
\end{description}

\begin{table}[H]
  \centering
  \caption{Resumo dos estudos de caso de auditoria.}
  \label{tab:audit-summary}
  \begin{tabularx}{\textwidth}{lclX}
    \toprule
    \textbf{Repositório} & \textbf{Nível} & \textbf{Licença} & \textbf{Uso recomendado} \\
    \midrule
    OdontoAI & G2 & MIT & Pesquisa \\
    DENTEX & G3 & Apache 2.0 & Produção clínica \\
    ToothSeg & G1+ & MIT & Aprendizado \\
    MeshSegNet & G2 & GPLv3 & Pesquisa 3D \\
    MLUA & G1 & Nenhuma (!) & \textbf{Não usar} \\
    DentalMIM & G2+ & MIT & Treinamento/Pesquisa \\
    \bottomrule
  \end{tabularx}
\end{table}

\section{Como Transformar um Repositório em Notebook Colab}

Uma habilidade prática essencial é converter o código de um repositório
GitHub em um notebook Google Colab executável:

\begin{enumerate}
  \item \textbf{Clone o repositório no Colab}:
\begin{lstlisting}[language=bash]
!git clone https://github.com/usuario/repo.git
%cd repo
\end{lstlisting}

  \item \textbf{Instale as dependências}:
\begin{lstlisting}[language=bash]
!pip install -r requirements.txt
\end{lstlisting}

  \item \textbf{Extraia as funções principais} para células do notebook,
    adicionando markdown explicativo entre elas.

  \item \textbf{Substitua paths hardcoded} por variáveis configuráveis
    ou uploads interativos (\texttt{files.upload()}).

  \item \textbf{Adicione visualizações} com matplotlib para cada etapa
    do pipeline.

  \item \textbf{Adicione testes TDD} usando \texttt{odontoia.tdd} para
    as funções críticas.
\end{enumerate}

\begin{pratica}{Auditoria de um repositório real}
  Objetivo: Escolher um repositório de IA odontológica no GitHub, aplicar
  o protocolo de auditoria G1-G3 e produzir um relatório.
\end{pratica}

\codigo{codigos/cap16/audit_github_repo.py}

\teste{O script de auditoria avalia automaticamente 15 critérios
distribuídos nos níveis G1, G2 e G3, atribuindo um nível de
confiabilidade e gerando recomendações específicas para cada repositório.}

\section{TDD Sobre Código de Terceiros: Um Exemplo}

Suponha que você encontrou um repositório G2 promissor com uma função
\texttt{detect\_caries\_bitewing()}, mas sem testes. Antes de usá-la:

\begin{lstlisting}[language=Python, caption=TDD sobre código de terceiros — etapa RED]
# test_audit_caries_detector.py
import numpy as np
from odontoia.tdd.test_helpers import create_test_radiograph

def test_detect_caries_output_format():
    """Verifica que a saída tem o formato documentado."""
    data = create_test_radiograph((224, 224), n_classes=2)
    img = data.images[0]

    # Importa a função do repositório auditado
    from third_party_caries import detect_caries_bitewing
    result = detect_caries_bitewing(img)

    # Contrato documentado no README:
    # result deve ser dict com 'has_caries', 'confidence', 'icdas'
    assert isinstance(result, dict)
    assert 'has_caries' in result
    assert isinstance(result['has_caries'], bool)
    assert 'confidence' in result
    assert 0.0 <= result['confidence'] <= 1.0
    print("[ICON] Saída segue contrato documentado")

def test_detect_caries_edge_cases():
    """Verifica comportamento em casos extremos."""
    from third_party_caries import detect_caries_bitewing

    # Imagem totalmente preta (subexposta)
    black_img = np.zeros((224, 224), dtype=np.uint8)
    result = detect_caries_bitewing(black_img)
    assert result['confidence'] < 0.5  # Não deve ter alta confiança

    # Imagem totalmente branca (superexposta)
    white_img = np.ones((224, 224), dtype=np.uint8) * 255
    result = detect_caries_bitewing(white_img)
    assert result['confidence'] < 0.5

    print("[ICON] Casos extremos tratados adequadamente")
\end{lstlisting}

\testefalha{Executando TDD sobre o repositório de terceiros: o teste
\texttt{test\_detect\_caries\_edge\_cases} falhou — a função original
retornava \texttt{confidence=0.98} para uma imagem totalmente preta,
indicando que o modelo não foi treinado com augmentação de brilho
adequada.}

\destaque{O TDD sobre código de terceiros é a \textbf{última linha de
defesa} antes de confiar em software externo para tarefas clínicas.
Nunca presuma que o código funciona como documentado — verifique.}

\section{Síntese do Capítulo}

Este capítulo forneceu um framework completo para auditar repositórios
de IA odontológica no GitHub:

\begin{itemize}
  \item \textbf{Busca avançada} com filtros por linguagem, stars, licença
    e data de atualização.
  \item \textbf{Níveis G1--G3} de confiabilidade: básico (exploratório),
    intermediário (pesquisa) e avançado (produção clínica).
  \item \textbf{Protocolo de auditoria em 5 etapas}: licença, artigo/DOI,
    dataset, execução/testes, TDD sobre terceiros.
  \item \textbf{Estudos de caso} de 6 repositórios reais, do G1 (MLUA —
    inutilizável) ao G3 (DENTEX — pronto para clínica).
  \item \textbf{TDD sobre código de terceiros} como prática de verificação
    independente antes da adoção.
\end{itemize}

A auditoria de repositórios não é burocracia: é \textbf{autodefesa
científica e clínica}.

\section{Além do GitHub: Outras Fontes de Código Odontológico}

\subsection{GitLab, Bitbucket e Repositórios Institucionais}

Embora o GitHub concentre a maioria dos projetos open source, repositórios
relevantes também estão em:

\begin{itemize}
  \item \textbf{GitLab}: preferido por instituições europeias (GDPR compliance).
  \item \textbf{Bitbucket}: comum em grupos de pesquisa que usam o ecossistema Atlassian.
  \item \textbf{Repositórios institucionais}: universidades frequentemente hospedam
    GitLab auto-hospedado (\textit{self-hosted}) com código de teses e dissertações.
  \item \textbf{Open Science Framework (OSF)}: mais voltado para dados e pré-registros,
    mas ocasionalmente contém código vinculado a artigos.
\end{itemize}

\subsection{Papers With Code e Model Zoo}

\begin{description}
  \item[\textbf{Papers With Code} (\url{https://paperswithcode.com})] é um
    diretório que vincula artigos científicos a implementações de código.
    Buscar por ``dental'', ``caries'' ou ``oral cancer'' frequentemente
    revela implementações não encontradas em buscas diretas no GitHub.

  \item[\textbf{Model Zoo}] repositórios como TensorFlow Hub, PyTorch Hub e
    ONNX Model Zoo oferecem modelos pré-treinados que podem ser
    reaproveitados. Um modelo de segmentação treinado em imagens médicas
    gerais pode ser fine-tunado para dentes.
\end{description}

\subsection{Google Dataset Search e Kaggle}

Datasets odontológicos públicos podem ser encontrados via:

\begin{itemize}
  \item \textbf{Google Dataset Search}: indexa datasets estruturados com
    metadados schema.org.
  \item \textbf{Kaggle}: abriga competições e datasets odontológicos, como
    o ``Dental Panoramic X-ray Dataset'' e ``Oral Cancer Images Dataset''.
    Verifique a licença (muitos datasets do Kaggle são CC0 ou CC BY).
\end{itemize}

\section{Ferramentas Automatizadas de Auditoria}

\subsection{Scorecard do OpenSSF}

A \textit{Open Source Security Foundation} (OpenSSF) mantém o
\textbf{Scorecard}, uma ferramenta automatizada que avalia a postura de
segurança de repositórios:

\begin{lstlisting}[language=bash]
# Instala e executa Scorecard
go install github.com/ossf/scorecard/v5@latest
scorecard --repo=github.com/odontoai/classification --format=json
\end{lstlisting}

O Scorecard verifica automaticamente:
\begin{itemize}
  \item Presença de licença e vulnerabilidades conhecidas.
  \item Assinatura de commits e proteção de branches.
  \item Frequência de atualizações e manutenção ativa.
  \item Uso de ferramentas de análise estática (SAST).
  \item Política de divulgação de vulnerabilidades.
\end{itemize}

\subsection{Reproducibility Check com REANA}

O \textbf{REANA} (REusable ANAlyses) é uma plataforma do CERN para
verificação de reprodutibilidade. Você pode submeter um repositório
e o REANA tentará reexecutar o pipeline de análise em um ambiente
containerizado.

\subsection{Ferramentas Específicas para IA}

\begin{itemize}
  \item \textbf{Model Cards Toolkit}: biblioteca do Google para criar e
    validar Model Cards.
  \item \textbf{TFX / ML Metadata}: rastreia a proveniência de cada
    artefato no pipeline de ML (dados, transformações, modelo).
  \item \textbf{DVC (Data Version Control)}: versiona datasets grandes
    junto com o código, essencial para reprodutibilidade.
\end{itemize}

\section{Red Flags: Sinais de Alerta em Repositórios}

Após auditar centenas de repositórios odontológicos, identificamos
padrões que devem acionar alarmes imediatos:

\begin{enumerate}
  \item \textbf{Repositório sem README ou com README vazio}: indica
    projeto abandonado ou nunca finalizado.

  \item \textbf{Autor anônimo ou sem afiliação institucional}: sem
    responsabilidade acadêmica, o código não tem lastro científico.

  \item \textbf{Números ``perfeitos demais''}: acurácia de 99,9\% em
    dataset de 50 imagens é virtualmente impossível e sugere overfitting
    extremo, data leakage ou manipulação.

  \item \textbf{Dependência de dataset proprietário sem descrição}:
    se você não sabe com quais dados o modelo foi treinado, não pode
    avaliar viés, validade ou generalização.

  \item \textbf{Código que só roda no ambiente do autor}: hardcoding de
    paths absolutos (\texttt{C:/Users/Fulano/Documents/dataset/}),
    dependências não listadas, scripts bash específicos de SO.

  \item \textbf{Único commit com todo o código}: \textit{``commit inicial:
    tudo''} com 50 arquivos sugere que o código foi despejado de uma só
    vez, sem histórico de desenvolvimento. Impossível auditar a evolução.

  \item \textbf{Issues ignoradas por mais de 6 meses}: comunidade faz
    perguntas, ninguém responde — projeto efetivamente abandonado.
\end{enumerate}

\destaque{Se um repositório apresenta 3 ou mais \textit{red flags},
considere-o \textbf{não auditável} e busque alternativas. O tempo gasto
tentando fazer código quebrado funcionar raramente compensa.}

\section{Estudo de Caso Expandido: Transformando G1 em G3}

Para ilustrar o valor prático da auditoria, documentamos o processo de
elevar um repositório odontológico real do nível G1 ao G3:

\begin{description}
  \item[Repositório original] \texttt{dental-caries-cnn} (G1)
  \item[Estado inicial] Código funcional, artigo no arXiv, sem testes,
    sem licença, sem CI/CD, sem Model Card.
\end{description}

\textbf{Intervenções realizadas:}

\begin{enumerate}
  \item \textbf{Fork e licença}: adicionada licença MIT após contato com
    o autor original.
  \item \textbf{Ambiente}: criado \texttt{environment.yml} com versões exatas.
  \item \textbf{Testes TDD}: adicionados 47 testes unitários usando
    \texttt{odontoia.tdd.test\_helpers}, cobrindo 82\% do código.
  \item \textbf{CI/CD}: GitHub Actions configurado para executar pytest
    e SpecVerifier a cada push.
  \item \textbf{Model Card}: documentação completa seguindo o template
    de Mitchell et al. (2019).
  \item \textbf{Docker}: criado Dockerfile para ambiente reproduzível.
  \item \textbf{Hugging Face}: modelo publicado com DOI via Zenodo.
\end{enumerate}

\textbf{Resultado}: O repositório passou de G1 (exploratório) para G3
(produção) em aproximadamente 20 horas de trabalho, tornando-se utilizável
para pesquisa clínica.

\teste{Após as intervenções, o repositório elevado a G3 passou em todos os
15 critérios de auditoria, com cobertura de testes de 82\% e CI/CD
verificando cada commit automaticamente.}

\section{Síntese do Capítulo}

Este capítulo forneceu um framework completo para auditar repositórios
de IA odontológica:

\begin{itemize}
  \item \textbf{Busca avançada} com filtros por linguagem, stars, licença
    e data de atualização — incluindo fontes além do GitHub (GitLab, OSF,
    Papers With Code, Kaggle).

  \item \textbf{Níveis G1--G3} de confiabilidade: básico (exploratório),
    intermediário (pesquisa) e avançado (produção clínica), com checklists
    detalhados para cada nível.

  \item \textbf{Protocolo de auditoria em 5 etapas}: licença, artigo/DOI,
    dataset, execução/testes, TDD sobre terceiros.

  \item \textbf{Estudos de caso} de 6 repositórios reais, do G1 (MLUA —
    inutilizável) ao G3 (DENTEX — pronto para clínica), incluindo o
    processo completo de elevação de G1 a G3.

  \item \textbf{Red flags}: 7 sinais de alerta que, se presentes, sugerem
    que o repositório não deve ser utilizado.

  \item \textbf{TDD sobre código de terceiros} como prática de verificação
    independente antes da adoção.

  \item \textbf{Ferramentas automatizadas}: OpenSSF Scorecard, REANA,
    Model Cards Toolkit, DVC.
\end{itemize}

A auditoria de repositórios não é burocracia: é \textbf{autodefesa
científica e clínica}. No próximo capítulo, abordaremos como garantir
que \textbf{seu próprio} código seja auditável, reproduzível e
corretamente licenciado.
